%% ============================================================
%%  PLANTILLA DE APUNTES DE MATEMÁTICA
%%  Compilar con:  pdflatex plantilla_apuntes.tex   (2 veces)
%%  Para el índice analítico: makeindex plantilla_apuntes
%%  y luego pdflatex otra vez. (Con --shell-escape se hace solo.)
%% ============================================================
\documentclass[11pt,a4paper,oneside]{report}

%% ---------- Idioma y codificación ----------
\usepackage[T1]{fontenc}
\usepackage[utf8]{inputenc}
\usepackage{lmodern}
\usepackage[spanish,es-noshorthands]{babel}
\usepackage{microtype}

%% ---------- Página y párrafos ----------
\usepackage[margin=2.5cm]{geometry}
\usepackage{parskip}            % separación entre párrafos sin sangría
\usepackage{setspace}
\setstretch{1.1}
\usepackage{fancyhdr}
\pagestyle{fancy}
\fancyhf{}
\fancyhead[L]{\small\nouppercase{\leftmark}}
\fancyhead[R]{\small\thepage}
\renewcommand{\headrulewidth}{0.4pt}

%% ---------- Matemática ----------
\usepackage{amsmath,amssymb,amsthm,mathtools}
\usepackage{bm}

%% ---------- Figuras, tablas, colores ----------
\usepackage{graphicx}
\graphicspath{{figuras/}}       % pon tus imágenes en la carpeta "figuras"
\usepackage{xcolor}
\usepackage{booktabs}           % tablas bonitas (\toprule, \midrule)
\usepackage{array}
\usepackage{tabularx}
\usepackage{longtable}
\usepackage{caption}
\usepackage{float}              % permite [H] para fijar posición
\usepackage{mwe}                % solo para imágenes de ejemplo; bórralo luego
\usepackage{tikz}

%% ---------- Código ----------
\usepackage{listings}
\lstset{
  basicstyle=\ttfamily\small,
  keywordstyle=\color{blue!70!black}\bfseries,
  commentstyle=\color{green!40!black}\itshape,
  stringstyle=\color{orange!80!black},
  numbers=left, numberstyle=\tiny\color{gray}, numbersep=8pt,
  backgroundcolor=\color{gray!8},
  frame=single, rulecolor=\color{gray!40},
  breaklines=true, showstringspaces=false, tabsize=4,
  literate={á}{{\'a}}1 {é}{{\'e}}1 {í}{{\'i}}1 {ó}{{\'o}}1 {ú}{{\'u}}1
           {Á}{{\'A}}1 {É}{{\'E}}1 {Í}{{\'I}}1 {Ó}{{\'O}}1 {Ú}{{\'U}}1
           {ñ}{{\~n}}1 {Ñ}{{\~N}}1
}

%% ---------- Índices ----------
\usepackage{imakeidx}
\makeindex[title=Índice analítico, intoc]

%% ---------- Cajas de colores ----------
\usepackage[most]{tcolorbox}

%% ---------- Hipervínculos y referencias (SIEMPRE al final) ----------
\usepackage[colorlinks=true,linkcolor=blue!60!black,citecolor=green!40!black,
            urlcolor=blue!60!black]{hyperref}
\usepackage[spanish,nameinlink]{cleveref}

%% ============================================================
%%  ENTORNOS: definición, teorema, demostración, notas...
%%  Todos comparten UN solo contador, dentro de cada capítulo.
%% ============================================================
\theoremstyle{definition}
\newtheorem{definicion}{Definición}[chapter]
\newtheorem{ejemplo}[definicion]{Ejemplo}
\newtheorem{ejercicio}[definicion]{Ejercicio}

\theoremstyle{plain}
\newtheorem{teorema}[definicion]{Teorema}
\newtheorem{proposicion}[definicion]{Proposición}
\newtheorem{lema}[definicion]{Lema}
\newtheorem{corolario}[definicion]{Corolario}

\theoremstyle{remark}
\newtheorem*{observacion}{Observación}
\newtheorem*{notademo}{Nota sobre la demostración}
\newtheorem*{idea}{Idea clave}

\renewcommand{\proofname}{Demostración}
\renewcommand{\qedsymbol}{$\blacksquare$}

% Nombres para cleveref
\crefname{definicion}{definición}{definiciones}
\Crefname{definicion}{Definición}{Definiciones}
\crefname{teorema}{teorema}{teoremas}
\Crefname{teorema}{Teorema}{Teoremas}
\crefname{proposicion}{proposición}{proposiciones}
\Crefname{proposicion}{Proposición}{Proposiciones}
\crefname{lema}{lema}{lemas}
\Crefname{lema}{Lema}{Lemas}
\crefname{corolario}{corolario}{corolarios}
\Crefname{corolario}{Corolario}{Corolarios}
\crefname{ejemplo}{ejemplo}{ejemplos}
\Crefname{ejemplo}{Ejemplo}{Ejemplos}
\crefname{ejercicio}{ejercicio}{ejercicios}
\Crefname{ejercicio}{Ejercicio}{Ejercicios}
\crefname{equation}{ecuación}{ecuaciones}
\Crefname{equation}{Ecuación}{Ecuaciones}
\crefname{figure}{figura}{figuras}
\Crefname{figure}{Figura}{Figuras}
\crefname{table}{tabla}{tablas}
\Crefname{table}{Tabla}{Tablas}
\crefname{chapter}{capítulo}{capítulos}
\Crefname{chapter}{Capítulo}{Capítulos}
\crefname{section}{sección}{secciones}
\Crefname{section}{Sección}{Secciones}

% Estilo visual: barra lateral de color + fondo suave
\tcbset{
  cajabase/.style={enhanced, breakable, sharp corners, boxrule=0pt,
    frame hidden, left=8pt, right=6pt, top=4pt, bottom=4pt,
    before skip=10pt, after skip=10pt}
}
\tcolorboxenvironment{definicion}{cajabase,
  colback=blue!6,   borderline west={3pt}{0pt}{blue!65!black}}
\tcolorboxenvironment{teorema}{cajabase,
  colback=red!6,    borderline west={3pt}{0pt}{red!70!black}}
\tcolorboxenvironment{proposicion}{cajabase,
  colback=orange!7, borderline west={3pt}{0pt}{orange!80!black}}
\tcolorboxenvironment{lema}{cajabase,
  colback=orange!7, borderline west={3pt}{0pt}{orange!80!black}}
\tcolorboxenvironment{corolario}{cajabase,
  colback=orange!7, borderline west={3pt}{0pt}{orange!80!black}}
\tcolorboxenvironment{ejemplo}{cajabase,
  colback=green!5,  borderline west={3pt}{0pt}{green!45!black}}
\tcolorboxenvironment{ejercicio}{cajabase,
  colback=green!5,  borderline west={3pt}{0pt}{green!45!black}}
\tcolorboxenvironment{notademo}{cajabase,
  colback=yellow!14, borderline west={3pt}{0pt}{yellow!60!black}}
\tcolorboxenvironment{observacion}{cajabase,
  colback=gray!10,  borderline west={3pt}{0pt}{gray!70}}
\tcolorboxenvironment{idea}{cajabase,
  colback=purple!6, borderline west={3pt}{0pt}{purple!70}}

%% ============================================================
%%  ATAJOS MATEMÁTICOS
%% ============================================================
\newcommand{\N}{\mathbb{N}}
\newcommand{\Z}{\mathbb{Z}}
\newcommand{\Q}{\mathbb{Q}}
\newcommand{\R}{\mathbb{R}}
\newcommand{\C}{\mathbb{C}}
\newcommand{\eps}{\varepsilon}
\newcommand{\dd}{\mathrm{d}}
\newcommand{\dpar}[2]{\frac{\partial #1}{\partial #2}}
\DeclarePairedDelimiter{\abs}{\lvert}{\rvert}
\DeclarePairedDelimiter{\norm}{\lVert}{\rVert}
\DeclarePairedDelimiter{\pesc}{\langle}{\rangle}
\DeclarePairedDelimiter{\conj}{\{}{\}}
\DeclareMathOperator{\sen}{sen}
\DeclareMathOperator{\tg}{tg}
\DeclareMathOperator{\Img}{Im}
\DeclareMathOperator{\Ker}{Ker}
\DeclareMathOperator{\rg}{rg}

% Atajo para "término indexado": lo escribe en el texto Y lo añade al índice
\newcommand{\termino}[1]{\textbf{#1}\index{#1}}

%% ============================================================
%%  PORTADA
%% ============================================================
\title{Apuntes de series de tiempo}
\author{Práxedis Jiménez Ruvalcaba}
\date{\today}

\begin{document}

\maketitle
\tableofcontents

%% ============================================================
\part{Series de Tiempo (Shumway, Stoffer)}
%% ============================================================

\chapter{Introducción}
Estos apuntes no son de un curso que se haya tomado de manera formal, sino que
mayormente apuntes autodidáctas del libro Time Series and Its Applications: With R Examples, 
de Shumway y Stoffer.

\section{Características de una serie de tiempo}
La necesidad de las series de tiempo viene de que contiene datos que incumplen la
"estacionariedad", de la estadística usual, es decir, los datos dependen del momento en 
tiempo en el que fueron tomados, más aún, los datos normalmente pertenecen a una dupla
$(valor,tiempo)$  ímplicita. Resultando en que tengan una autocorrelación, incumpliendo también las 
hipótesis de independencia e idénticamente distribuidas. El acercamiento sistémico a la
solución de estos problemas son las series de tiempo.\\

Una recomendación escencial es la de realizar un escrutinio claro de la información presente en
la gráfica al comienzo de cualquier análisis o investigación de una serie de tiempo. Muchas veces,
a través del gráfico y el contexto, podemos seleccionar el método de análisis estadístico para 
condensar propiamente la información. Los enfoques particulares de esto último caen en dos categorías:
$domain$ $approach$ y $frequency$ $domain$ $approach$.\\

\textbf{Domain Approach:}\\
Se asume que la correlación entre las observaciones adyacentes; pueden expresar el valor presente con 
los valores pasados $pasado \rightarrow presente$. Modelar eventos futuros a través de información 
pasada rima en la estadística más clásica con las regresiones y es que precisamente este enfoque 
trabaja realizando regresiones lineales del presente al pasado, para generar información. 
De estás ideas surgen las \textbf{ARIMA} ($autoregresive$ $integrated$ $meaning$ $average$) 
para operar modelos con correlación en base al tiempo. A su vez existe un formato multivariado para 
ARIMA o el uso de funciones de transferencias para cumplir con estos mismos objetivos.\\
Estos modelos son multiplicativos, por lo que sus datos se asumen provenir de factores
respecto a un "ruido blanco" entrante.

\begin{observacion}
  El nombre de "ruido blanco" viene de la luz blanca porque sus ondas poseen un comportamiento$\dots$
  que no sé describir bien porque no soy de física y no llevé óptica, pero ahí lo dejo.
\end{observacion}

Existen modelos que provienen más de la suma de series que de un producto de factores, donde 
cada sumando es una cierta estructura de serie de tiempo específica. Esto facilita muchas 
consideraciones y procesos, pero también viene con sus complicaciones y limitaciones para ciertas
condiciones.\\
Las series de tiempo en memoria larga como GARCH también son un enfoque válido bajo el marco que 
discute el $domain$ $approach$.\\

\textbf{Frequency Domain:}\\
En lugar de un correlación de datos adyacentes donde esta medida depende de la distancia entre
puntos, este enfoque asume que la correlación es por una periodicidad, ya no depende de la 
distancia entre puntos sino de los puntos en si. Normalmente se describe como donde hay presente
una variación sinuidosidal.\\
En Análisis espectral, la partición de varios tipos de variación en una serie de tiempo se 
trabaja evalúando por separando las varianzas asociadas en la de periodicidad de interés.\\

\begin{definicion}[\termino{Serie de Tiempo}]
  Colección de variables aleatorias indexadas acorde al orden el que se realizaron las observaciones
  de los datos acorde a una medida de tiempo.\\
  Sea $x_i$ una variable aleatoria $\Rightarrow x_1, x_2, x_3,...$ es una ts donde $x_1$ es la variable aleatoria 
  correspondiente al valor observado en el primer momento, $x_2$ en el segundo,$...$\\
  La serie puede escribirse como $x_t$ o $\conj{x_t}$ usualmente.
\end{definicion}

$\conj{x_t}$ propiamente se refiere a un proceso estocástico donde se aume $t$ discreto, ya que 
obtener datos de manera continua y evaluarlo es algo prácitcamente inviable. Es por esta clase de
detalles que debe tomarse en cuenta para entender que la elección de intervalos en los cuales muestrear
los datos para la serie de tiempo, es escencial. De no ser seleccionado un apropiado intervalo caemos en
$aliasing$ que es cuando el muestreo es insuficiente y termina generando algo totalmente distinto en los datos,
es como el efecto Moire o las sierritas en los gráficos 3D.

\begin{ejemplo}[Ruido Blanco]\label{ej:ruido}
  Sea $w_i ~ N(0,\sigma^2) \forall i$ $\Rightarrow \conj{w_t}$ solo tiene datos aleatorios sin correlación
  alguna entre ellas, es lo que normalmente se entiende como ruido blanco. Esta variedad de series con 
  son de sumo interés y normalmente se describe al ruido con la antes mencionada distribución. 
\end{ejemplo}
  En varios ejemplos y ejercicios va a ser más concretamente con una normal estándar en lugar de varianza 
  arbitraria. Lo que más importa es la distribución normal con media cero.\\

Si las ts pudieran ser explicadas en términos del ruido blanco, podríamos utilizar la estadística usual. 
El detalle es que las ts no tienen este comportamiento usualmente o uno que nos permita de manera tan sencialla
trabajar de esta forma.

\begin{ejemplo}[Media Deslizante]
Consideremos $w_t$ del ejemplo anterior y realizemos una ponderación de valores:
\[ x_t=\frac{1}{3}(w_{t-1}+w_{t}+w_{t+1}) \]
\end{ejemplo}
Se puede notar que está tendrá un comportamiento similar a \cref{ej:ruido} pero suavizando mucho, especialmente
si se les compara en la misma escala. Es posible que si los gráficas por separado, sean parecidos pero algo recorridos,
la gracia es cuando están en la misma escala y el haber calculado una media a través de los valores se nota, literalmente
mediando los valores por los que va pasando.

\begin{ejemplo}[Autoregresión]
  Ahora consideremos
  \[x_t=x_{t_1}-0.90x_{t-2} +w_t \]

\end{ejemplo}
Como $x_t$ depende de dos casos previos, en realidad se está haciendo una regresión en base a los valores previos,
una predicción. Vale la pena revisar cómo se compara con \cref{ej:ruido} y con el anterior. En general, 
es bueno comparar todos los ejercicios previos. 

\begin{ejemplo}[Random Walk]
  \[ x_t=\delta + x_{t-1} + w_t\]
  Consideremos además $x_0=0$. Entonces $w_t$ es el ruido y $\delta$ es el $drift$. 
\end{ejemplo}
El $drift$ es la dirección de la caminata, hacía dónde se encaminan o con qué pendiente deberían ir 
variando los valores. Este tiempo de ts se utiliza mucho en análisis financieros, particularmente
para la simulación de variaciones de precio a través de Monte Carlo para determinar con intervalos
de confianza la manera en que puede ir cambiando el precio con el tiempo, aunque es cierto que 
el mismo problema usualmente es resuelto de manera anlítica con las ecuaciones Black-Scholes.\\
Notemos además que $x_t=\delta t + \sum_{j = 1}^{t}w_j$

\begin{proof}
  $x_t=\delta t + \sum_{j = 1}^{t}w_j$\\
  Resolvamos por inducción\\
  La proposición es válida para el primer valor admisible, recordemos $x_0=0$\\
  $\Rightarrow x_1=\delta + x_0+w_1=\delta + w_1$\\
  Consideremos que la proposición es válida para un valor arbitrario, digamos $k$, nos
  daría la hipótesis de inducción:
  \[x_k=\delta k + \sum_{j = 1}^{k}w_j\]
  Revisemos si la proposición es válida para el siguiente valor admisible, digamos $k+1$\\
  $x_{k+1}=\delta + x_k + w_{k+1}$\\
  $\quad \delta + \delta k + \sum_{j = 1}^{k}w_j + w_{k+1} $ por hipótesis de inducción\\
  $\quad \delta (k+1) + \sum_{j = 1}^{k+1}w_j$
\end{proof}  

\begin{ejemplo}[Señales en ruido]
\[x_t=2cos(\frac{2\pi t}{50}+0.6\pi)+w_t \]
En términos generales podemos describir una ts con forma sinuidoal como:
\[Acos(2\pi wt + \phi)\]
donde
\begin{itemize}
  \item $A=$ amplitud
  \item $w=$ frecuencia de la oscilación
  \item $\phi=$ cambio de fase 
\end{itemize}
\end{ejemplo}
El punto es observar cuando hay una variación de período consistente con un ruido aleatorio añadido\\
Casi todos los ejemplos se trataron de modelos simples que son aquellos que pueden expresarse
como $x_t=s_t+v_t$ donde $s_t$ son los datos de la serie y $v_t$ es tanto el ruido como si hay algo
correlacionado. 

\section{Autocorrelación}
La manera en que unos datos dependen de otros no orientan a determinar que si hay una medida en 
el que es particular interés fijarse son la covarianza y la correlación. \\
Una serie de tiempo son al final $n$ variables aleatorias en $n$ momentos $t_1, t_2,...,t_n$,
así que la distribución conjunta de estas nos daría toda la información que requisieramos para
analizar propiamente la serie de tiempo. En particular, la distribución conjunta evaluando
la probabilidad de los valores de la serie sean menores que las constantes $c_1,c_2,...,c_n$
\[F(c_1,c_2,...,c_n)=P(x_{t_1}\leq c_1,x_{t_2}\leq c_2,...,x_{t_n}\leq c_n)\]

El detalle es que una función así solo sería sencilla de obtener y evaluar si tratase de una
normal multivariada o cada entrada fuera iid normal porque entonces podríamos por teorema de 
probabilidad expresar la distribución de probabilidad conjunto como un producto de marginales
\[F(c_1,c_2,...,c_n)=\prod_{t=1}^{n}\phi(c_t)\] 
\[\phi(x)=\frac{1}{\sqrt{2\pi}}\int_{0}^{x}e^{-\frac{z^2}{2}} dz \]
Que es la acumulada de una normal.\\
Y una vez conseguido describir los datos de esta forma, el otro detalle es que no se puede gráficar.
Lo más que se puede hacer es graficar la distribución de un momento específico. Tal vez dos e incluso
ahí no tendría mucho sentido ni aporte, admeás de que lo hace parecer como si quisiera uno forzar
los usos y costumbres de la estadística usual.

\begin{definicion}[\termino{Función Media}]
  \[ \mu_{x_t}=E(x_t)=\int_{-\infty}^{\infty}xf_t(x)dx\]
\end{definicion}
$E$ es el operador de valor esperado que vuelve de probabilidad con el mismo significado como
primer momento de centralidad. Solo que es la centralidad en el momento $t$, se analiza hasta 
lo que en teoría ha causado o repercutido dicho momento.

\begin{ejemplo}\label{ej:media RW}
  Retomando los ejemplos de la sección anterior, empezando por \cref{ej:ruido}\\
  1)\\
  $E(w_t)=0$ porque $w_i~N(0,\sigma^2) \forall i$\\
  2)\\
  $E(\frac{1}{3}(w_{t-1}+w_{t}+w_{t+1}))=\frac{1}{3}(E(w_{t-1})+E(w_t)+E(w_{t+1}))$\\
  $\quad = \frac{1}{3}(0)=0$ por 1)\\
  3)\\
  $E(\delta t + \sum_{j = 1}^{t}w_j)=E(\delta t)+\sum_{j = 1}^{t}E(w_j)$ por linealidad de $E$\\
  $\quad \delta t +0 =\delta t$ El drift por el tiempo es la media, lo cuál es una recta.
\end{ejemplo}
Notése que al ser una recta, esa última media no está acotada por arriba, por lo que podría
no ser finita, algo que impediría la existencia de varianza en la distribución. 

\begin{definicion}[\termino{Función de Autocovarianza}]
  \[\gamma_x(s,t)=E[(x_s-\mu_s)(x_t-\mu_t)] \]
\end{definicion}
Se mantiene la autocovarianza bajo la idea de del segundo momento de centralidad, o un producto para constuir
un área esperada alrededor del primer momento. En particular, mide la dependencia lineal entre dos momentos.\\
Una serie suave va a tomar valores grandes de $\gamma$ en momentos muy separados entre si, mientras que una serie
de mcuhos cambios bruscos y repentienos va a tomar valores muy cercanos a cero o suficientemente bajos. De hecho, el concepto
de autocovarianza se va reformular en base a la "estacionariedad" y precisamente, la distancia entre puntos marcará
la autocovarianza de la serie de tiempo.\\
Recordemos $E(xy)$ de normal en teoría de probabilidad mide cuánto varían conjuntamente de media $X$ y $Y$. Es una media
de varianza para la distribución conjunta.\\
\begin{observacion}
  Si $\gamma_x(s,t)=0$\\
  Implica una falta de dependencia lineal, pero no una ausencia total de factores de dependencia.
  La igualdad estricta a 0 pues signifca una ortogonalidad. A menos de que ambos operan iid normales,
  entonces el 0 significa independencia.
\end{observacion}

\begin{ejemplo}
  Tomando $w_t~N(0,\sigma^2)$ de \cref{ej:ruido}\\
  $\sigma_w(s,t)=E[(w_s-\mu_s)(w_t-\mu_t)]$\\
  $E[(w_s)(w_t)]$\\
  Si $s=t$ queda claro que $E[(w_s)(w_t)]=\sigma^2$ por teoría de probabilidad.
\end{ejemplo}
Cuando $s\not =t$, al ser pares al azar (no hay correlación) y por otro resultado de probabilidad, 
como las medias de los calores del ruido blanco no varían, entonces $E[(w_s)(w_t)]=0$.

\begin{ejemplo}
  Consideremos $s=t$ para el caso de la media deslizante\\
  $\gamma(s,t)=\frac{1}{9}E[(w_{t-1}+w_{t}+w_{t+1})(w_{t-1}+w_{t}+w_{t+1})]$\\
  Por el ejemplo anterior, cuando $w_s\not = w_t\Rightarrow E(w_ww_t)=0$, así que tendríamos
  $\gamma(s,t)=\frac{1}{9}[E(w_{t-1}w_{t-1})+E(w_{t}w_t)+E(w_{t+1}w_{t+1})]=\frac{1}{3}$ si consideramos $\sigma^2=1$\\
  De manera similar obtendríamos que para $s=t+1 \Rightarrow \gamma(s,t)=\frac{2}{9}$, $s=t+2 \Rightarrow \gamma(s,t)=\frac{1}{9}$,
  y para $s=t+n | n>2 \Rightarrow \gamma(s,t)=0$.
\end{ejemplo}
Cabe destacar que los resultados anteriores son idénticos si se toman las cantidad negativas con respecto a $t$.
Esto ya nos invita a pensar en la autocovarianza como algo que varía no por los puntos en si, si no por la distancia
que poseen los puntos que evalúamos entre sí, siendo los valores de la distancia invariantes dado que $s=t+1$ es lo 
mismo que $s=t-1$ en su autocovarianza para esto caso. Pensar en esto es escencial para cuando se redefinan estos conceptos en función
de estacionalidad, que va a ser precisamente el que todo bueda describirse en términos del lag una parte escencial de la definición.

\begin{ejemplo}
  Es de sumo interés revisar el caso del Random Walk, concretamente \cref{ej:media RW}, donde vimos que
  $E(x_t)=\delta t$. Veamos que \\
  $\sigma(s,t)=E[((\delta t + \sum^{t}w)-\delta t )((\delta s + \sum^{s}w)-\delta s )=E[(\sum^{t}w)(\sum^{s}w)]$
  Nuevamente, por el primer ejemplo entendemos que $E(w_jw_i)\not = 0 \iff i=j$ y esto solo es en pares de mismo índice 
  cuando las sumatorias coincidan, es decir, $\gamma(s,t)=min(s,t)\sigma^{2}$
\end{ejemplo}
A diferencia del anterior donde la covarianza está dada por la distancia entre los puntos independeinte de los puntos en sí, 
aquí la autocovarianza depende explícitamente de los puntos, el momento mínimo es un multiplicador directo de la autocovarianza.
Entre más cercano al inicio sea uno de los puntos, menor será la autocovarianza, pero si ambos se alejan crece sin cota, en particular;
$\lim_{t\rightarrow \infty }\gamma (t,t)=\lim_{t\rightarrow \infty }t\sigma^{2}=\infty$, que va de la mano con la condición de que
solo hay varianza si existe media definida y ya habíamos visto en \cref{ej:media RW}  que la media se indeterminaba para 
puntos suficientemente lejanos en cuanto a tiempo. No hay una cota con respecto al tiempo dado el incremento (o decremento) paulatino
constante lineal introducido por el término de $drift$.

\begin{definicion}[\termino{Función de Autocorrelación (ACF)}]
  \[ \rho (s,t)=\frac{\gamma (s,t)}{\sqrt{\gamma(s,s)\gamma(t,t)}} \]
\end{definicion}
La garantía de que el valor de la correlación se comporte como normalmente la entendemos está dada por Cauchy-Schwarz,
$|\gamma(s,t)|^2\leq \gamma(s,s)\gamma(t,t)$ implica $-1\leq p(s,t)\leq 1$. Este es el caso para todas las definiciones
de los distintos tipos de ACF a menos de que se especifique. \\
La ACF nos permite medir que tan predicible es el momento $t$ usando el $s$. Nos invita inmediatamente a pensar no solo 
en momentos, sino en otras series. Pero esto rquiere matizar también la autocovarianza.

\begin{definicion}[función de \termino{covarianza cruzada}]
  \[\gamma_{xy}(s,t)=E[(x_s-\mu_{xs})(y_t -\mu_{yt})]\]
\end{definicion}
Una medida para momentos distintos de series distintas. Tal vez algo muy general para que sea realmente práctico. Pero
ya se ajustará un poco con el siguiente concepto tras readaptar la ACF para comparación.

\begin{definicion}[Función de \termino{correlación cruzada} (CCF)]
    \[ \rho_{xy} (s,t)=\frac{\gamma_{xy} (s,t)}{\sqrt{\gamma_x(s,s)\gamma_y(t,t)}} \]
\end{definicion}

\section{Estacionariedad}

\begin{definicion}[Serie Estacionaria Estricta]
  Una ts para la que comportamiento probabílistico de cualesquier colección de valores $\conj{x_1,x_2,...,x_n}$
  es idéntico a la colección con tiempo reajustado $\conj{x_{1+h},x_{2+h},...,x_{n+h}}$. Es decir,
  \[ P\conj{x_1\leq c_1,..., x_n \leq c_n}=P{x_{1+h}\leq c_1,..., x_{1+h} \leq c_n}\]
  Para cualesquiera $n$ momentos de tiempo, todos números $c_1,...,c_n$, y todos los cambios de tiempo $h=0,|1|,|2|,...$
\end{definicion}
Como todo debe tener la misma probabilidad en cualesquier momento, la distribución y por tanto parámetros como la media, 
también deben ser siempre iguales. Como la media entonces no varía nunca, se trata de una constante. A su vez, la autocovarianza
entre puntos debería ser siempre la misma, ya que el cambio de distancia no debe de cambiar. Es decir,
$\gamma(s,t)=\gamma(s+h,t+h)$ por cómo deben mantenerse las probabilidades. La autocovarianza entonces solo depende de la distancia
entre momentos y no de los momentos en si.\\
El problema es que esta definción es demasiado estricta, pero la útilidad de depender solo de la distancia es demasiado tentador, así
que la condición de que debe cumplir para cualequier $n$-colección se debilita a que funcione solo para un par. 

\begin{definicion}[\termino{Serie Estacionaria}]
  $x_t$ es un proceso de varianza finita tal que 
  \begin{enumerate}
    \item la función media, $\mu_t$ es una constante que no depende de $t$ 
    \item La función de covarianza $\gamma(s,t)$ depende meramente de $|s-t|$
  \end{enumerate}
  Estrictamente estacionario $\Rightarrow$ estacionario
\end{definicion}

Lo que nos da las funciones adaptadas de autocovarianza y autocorrelación
  \[\gamma_{}(h)=E[(x_{t+h}-\mu)(x_t -\mu_)]\]
  \[ \rho(h)=\frac{\gamma (t+h,t)}{\sqrt{\gamma(t+h,t+h)\gamma(t,t)}}=\frac{\gamma(h)}{\gamma(0)} \]
Expresiones las cuáles facilitan la interpretación y cálculos.
\begin{observacion}
  Por Cauchy-Schwarz que nos garantizaba el comportamiento apropiada de $\rho$, también garantiza que 
  $\gamma (0)$ sea la varianza de la ts ya que $|\gamma(h)|\leq \gamma(0)$. Esto solo en el caso de 
  ser estacionaria.
\end{observacion}
De manera similar podemos definir los casos de comparación o de "cruce"
  \[\gamma_{xy}(h)=E[(x_{t+h}-\mu)(y_t -\mu_)]\]
  Como una función de covarianza cruzada que si se cumple que solo depende de $h$, podemos 
  decir que las series $x_t$ y $y_t$ son conjuntamente estacionarias. También deben cumplir 
  que ambas sean estacionarias.
  \[ \rho_{xy}(h)=\frac{\gamma_{xy}(h)}{\sqrt{\gamma_x(0)\gamma(0)}}\]

\section{Ejercicios}
\begin{ejercicio}
Considere la ts $x_t=s_t+w_t$ donde $w_t$ es ruido blanco gaussiano con media 0 y desviación 1. 
Simule y grafique los valores para $n=200$. Las formas de $s_t$ son:\\

a)\\
$s_t=\left\{
  \begin{array}{l}
    0 \quad 1 \leq t \leq 100 \\
    10e^{\frac{-(t-100)}{20}}cos(\frac{2 \pi t}{4}) \quad 101\leq t \leq 200
  \end{array}
\right.$\\

b)\\
$s_t=\left\{
  \begin{array}{l}
    0 \quad 1 \leq t \leq 100 \\
    10e^{\frac{-(t-100)}{200}}cos(\frac{2 \pi t}{4}) \quad 101\leq t \leq 200
  \end{array}
\right.$\\

c)\\
Compara y analiza ambas, además de revisar cómo se comportan en especifico los moduladores
de ambas ts. Considera solo $t=1,...,100$

\end{ejercicio}

\begin{lstlisting}[language=R, caption={Solución}]
zeros<-rep(0,100)
w<-rnorm(200,0,1)
t<-101:200
#a)
s<-10*exp(-(t-100)/20)*cos((2*pi*t)/4)
x<-append(zeros,s) #Me recuerda a la curva del tópologo
plot.ts(x+w)

#b)
s<-10*exp(-(t-100)/200)*cos((2*pi*t)/4)
x<-append(zeros,s) 
plot.ts(x+w)

#c)
#Moduladores de señal a)
t<-1:100
modula<-exp(-(t-100)/20)
plot.ts(modula) 
#moduladores de b)
modula<-exp(-(t-100)/200)
plot.ts(modula)


\end{lstlisting}
Reducir el exponente del producto exponencial reduce el ritmo de encoge de los picos de los períodos en b). \\
La primera es bastante más similar a los gráficos del terremoto, aunque la parte media es mucho más
ordenada que la del terremoto .La segunda se parece más a la explosión en el desorden inicial y mayor 
patrón al final, aunque no captura el descenso repentino que la anterior sí tenía, aquí los picos
se mantienen de forma mucho más similar a la del terremoto.\\
Exponencial negativa de pendiente pronunciada como el descenso del gráfico, 
desciende a su décima parte en menos de la mitad del tiempo.Parece recta del poco cambio que hay, 
es casi que constante en su descenso y no desciende ni la mitad de su valor en el total del tiempo

%% ============================================================
%%  ÍNDICE ANALÍTICO
%% ============================================================
\printindex

\end{document}
